\subsection{直和的对称代数。自由模的对称代数。分次模的对称代数}

设 A 为一个交换环， \(M = \bigoplus_{\lambda \in L} M_{\lambda}\) 为一族 A-模的直和， \(j_{\lambda}: M_{\lambda} \to M\) 为典范单射；我们导出一个 A-代数同态 \(\mathsf{S}(j_{\lambda}): \mathsf{S}(\mathsf{M}_{\lambda}) \to \mathsf{S}(\mathsf{M})\) 。由于 \(\mathsf{S}(\mathsf{M})\) 是交换的，§ 4，第 5 号，命题 8 可应用于诸同态 \(\mathsf{S}(j_{\lambda})\) ，因此存在唯一的代数同态

\[
g \colon \bigotimes_ {\lambda \in L} \mathbf {S} (\mathbf {M} _ {\lambda}) \rightarrow \mathbf {S} (\mathbf {M})\tag{10}
\]

（也记作 \(g_{\mathbf{M}}\) ）使得 \(\mathbf{S}(j_{\lambda}) = g \circ f_{\lambda}\) （对所有 \(\lambda \in \mathbf{L}\) ），其中

\[
f _ {\lambda} \colon \mathbf {S} (\mathbf {M} _ {\lambda}) \rightarrow \bigotimes_ {\lambda \in \mathbf {L}} \mathbf {S} (\mathbf {M} _ {\lambda})
\]

表示典范同态。

命题 9. 典范同态 \(g\) （公式 (10)）是一个分次代数同构（参见 § 4，第 8 号，注 1）。

为了证明 \(g\) 是双射的，我们定义一个代数同态

\begin{enumerate}
\def\labelenumi{(\arabic{enumi})}
\setcounter{enumi}{10}
\tightlist
\item
\end{enumerate}

\[
h \colon \mathbf {S} (\mathbf {M}) \rightarrow \bigotimes_ {\lambda \in L} \mathbf {S} (\mathbf {M} _ {\lambda})
\]

使得 \(g \circ h\) 和 \(h \circ g\) 分别是 \(\mathbf{S}(\mathbf{M})\) 和 \(\bigotimes_{\lambda \in \mathbf{L}} \mathbf{S}(\mathbf{M}_{\lambda})\) 上的恒等映射。对于每个 \(\lambda \in \mathbf{L}\) ，设 \(u_{\lambda}\) 为复合线性映射

\[
\mathbf {M} _ {\lambda} \xrightarrow {\phi_ {\mathrm{M} _ {\lambda}} ^ {\prime}} \mathbf {S} (\mathbf {M} _ {\lambda}) \xrightarrow {f _ {\lambda}} \bigotimes_ {\lambda \in L} \mathbf {S} (\mathbf {M} _ {\lambda}).
\]

存在唯一的 A-线性映射 \(u: \mathbf{M} \to \bigotimes_{\lambda \in L} \mathbf{S}(\mathbf{M}_{\lambda})\) ，使得 \(u \circ j_{\lambda} = u_{\lambda}\) （对所有 \(\lambda \in L\) ）。由于诸 \(\mathbf{S}(\mathbf{M}_{\lambda})\) 是交换的，它们的张量积也是交换的（§ 4，第 5 号），因此（第 1 号，命题 2）存在唯一的代数同态 \(h: \mathbf{S}(\mathbf{M}) \to \bigotimes_{\lambda \in L} \mathbf{S}(\mathbf{M}_{\lambda})\) ，使得 \(h \circ \phi_{\mathbf{M}}' = u\) ；另一方面，显然 \(u(\mathbf{M})\) 包含于分次代数 \(\bigotimes_{\lambda \in L} \mathbf{S}(\mathbf{M}_{\lambda})\) 的次数为 1 的元素所构成的子模中，因此 \(h\) 是一个分次代数同态。对于 \(x_{\lambda} \in \mathbf{M}_{\lambda}\) ，

\[
h \left(g \left(u _ {\lambda} \left(x _ {\lambda}\right)\right)\right) = h \left(g \left(f _ {\lambda} \left(\phi_ {M _ {\lambda}} ^ {\prime} \left(x _ {\lambda}\right)\right)\right)\right) = h \left(S \left(j _ {\lambda}\right) \left(\phi_ {M _ {\lambda}} ^ {\prime} \left(x _ {\lambda}\right)\right)\right) = h \left(\phi_ {M} ^ {\prime} \left(j _ {\lambda} \left(x _ {\lambda}\right)\right)\right) = u _ {\lambda} \left(x _ {\lambda}\right);
\]

由于诸 \(u_{\lambda}(x_{\lambda})\) 生成代数 \(\bigotimes_{\lambda \in L} S(M_{\lambda})\) （§ 4，第 5 号，命题 8）， \(h \circ g\) 当然是恒等映射。类似地，

\[
g \left(h \left(\phi_ {\mathrm{M}} ^ {\prime} \left(j _ {\lambda} \left(x _ {\lambda}\right)\right)\right)\right) = g \left(u _ {\lambda} \left(x _ {\lambda}\right)\right) = g \left(f _ {\lambda} \left(\phi_ {\mathrm{M} _ {\lambda}} ^ {\prime} \left(x _ {\lambda}\right)\right)\right) = \mathrm{S} \left(j _ {\lambda}\right) \left(\phi_ {\mathrm{M} _ {\lambda}} ^ {\prime} \left(x _ {\lambda}\right)\right) = \phi_ {\mathrm{M}} ^ {\prime} \left(j _ {\lambda} \left(x _ {\lambda}\right)\right)
\]

并且，由于诸元素 \(\phi_{\mathbf{M}}^{\prime}(j_{\lambda}(x_{\lambda}))\) 生成代数 \(\mathbf{S}(\mathbf{M})\) ， \(g \circ h\) 当然是恒等映射。

注记 (1) 设 \(\mathbf{N} = \bigoplus_{\lambda \in L} \mathbf{N}_{\lambda}\) 是另一族以 \(L\) 为指标集的 A-模的直和，并且对所有 \(\lambda \in L\) ，设 \(v_{\lambda}: \mathbf{M}_{\lambda} \to \mathbf{N}_{\lambda}\) 是一个 A-线性映射，由此存在一个 A-线性映射 \(v = \bigoplus_{\lambda} v_{\lambda}: \mathbf{M} \to \mathbf{N}\) (II，§1，第6号，命题6)。则图表

\[
\begin{array}{c} \bigotimes_ {\lambda \in L} S (M _ {\lambda}) \xrightarrow {g _ {M}} S (M) \\ \bigoplus_ {\lambda \in L} S (v _ {\lambda}) \Biggl \downarrow \qquad \qquad \qquad \qquad \qquad \qquad \qquad \qquad \qquad \Biggl \downarrow S (v) \\ \bigotimes_ {\lambda \in L} S (N _ {\lambda}) \xrightarrow {g _ {N}} S (N) \end{array}
\]

是交换的，这由定义直接得出（§ 4，第 5 号，命题 8 的推论）。

\(\bigotimes_{\lambda \in L} \mathbf{S}(\mathbf{M}_{\lambda})\) 中与 \(\mathbf{S}^n(\mathbf{M})\) 通过同构 \(g\) 等同起来的那个 A-子模可以更精确地描述。对于每个有限子集 \(J\) （ \(L\) 中的），我们记 \(\mathbf{E}_J = \bigotimes_{\lambda \in J} \mathbf{S}(\mathbf{M}_{\lambda})\) ，使得 \(\bigotimes_{\lambda \in L} \mathbf{S}(\mathbf{M}_{\lambda}) = \varinjlim \mathbf{E}_J\) （相对于有向集 \(\mathfrak{F}(\mathrm{L})\) ，即由 \(\mathbf{L}\) 的有限子集组成的集，根据定义，§ 4，第 5 号）。对于每个族 \(\nu = (n_{\lambda}) \in \mathbf{N}^{(\mathbf{L})}\) （因而具有有限支撑）满足 \(\sum_{\lambda \in \mathbf{L}} n_{\lambda} = n\) ，以及每个有限子集 \(\mathbf{J}\) （ \(\mathbf{L}\) 中的）——包含族 \(\nu\) 的支集——，我们记

\[
\mathbf {S} ^ {\mathrm{J}, \nu} (\mathbf {M}) = \bigotimes_ {\lambda \in J} \mathbf {S} ^ {n _ {\lambda}} (\mathbf {M} _ {\lambda})\tag{12}
\]

使得子模 \(E_{J,n}\) （由次数为 n 的、属于 \(E_{J}\) 的元素组成）是诸 \(S^{J,v}(M)\) 的直和，遍历所有支集包含于 J 且满足 \(\sum_{\lambda\in L}n_{\lambda}=n\) 的族 v（ \(\S4\) ，第7号，命题10和 \(\S4\) ，第8号）。作为约定，我们对支集不包含在 J 中的族 v 写 \(S^{J,v}(M)=\{0\}\) ；那么 \(E_{J,n}\) 也可以写成所有 \(S^{J,v}(M)\) 的直和，其中 v 遍历集合 \(H_{n}\) ，即所有族 \(v=(n_{\lambda})_{\lambda\in L}\) 中满足 \(\sum_{\lambda\in L}n_{\lambda}=n\) 的那些族组成的集合。由于 \(S^{0}(M_{\lambda})\) 与 A 等同，显然还有：对于 L 的两个有限子集 \(J\subset J'\) 以及一个支集包含于 J 中的族 v，典范映射 \(S^{J,v}(M)\to S^{J',v}(M)\) （典范映射 \(E_{J}\to E_{J'}\) 在 \(S^{J,v}(M)\) 上的限制）是双射。如果对所有 \(v\in H_{n}\) 我们写

\[
\mathbf {S} ^ {\nu} (\mathbf {M}) = \varinjlim \mathbf {S} ^ {\mathrm{J}, \nu} (\mathbf {M})\tag{13}
\]

由此可见，结合II，§6，第2号，命题5：

推论。A-模 \(\mathbf{S}^n (\mathbf{M})\) 是 \(\bigotimes_{\lambda \in \mathbf{L}}\mathbf{S}(\mathbf{M}_{\lambda})\) 的这样一个子模在同构 (10) 下的像：该子模是诸子模 \(\mathbf{S}^{\nu}(\mathbf{M})\) 的直和，遍历所有族 \(\nu = (n_{\lambda})\in \mathbf{N}^{(L)}\) （满足 \(\sum_{\lambda \in \mathbf{L}}n_{\lambda} = n\) ）；若 \(\mathbf{J}\) 是 \(\nu\) 的支集，则 \(\mathbf{S}^{\nu}(\mathbf{M})\) 典范同构于 \(\bigotimes_{\lambda \in \mathbf{J}}\mathrm{S}^{n_{\lambda}}(\mathbf{M}_{\lambda})\) 。

一般地， \(\mathbf{S}^{\nu}(\mathbf{M}),\bigotimes_{\lambda \in J}\mathbf{S}^{n_{\lambda}}(\mathbf{M}_{\lambda})\) 以及它们在 \(\mathbf{S}^n (\mathbf{M})\) 中的像被等同起来。

定理 1. 设 A 是一个交换环，M 是一个以 \((e_{\lambda})_{\lambda \in L}\) 为基的自由 A-模。对于每个具有有限支撑的映射 \(\alpha: \mathbf{L} \to \mathbf{N}\) ，我们记

\[
e ^ {\alpha} = \prod_ {\lambda \in L} e _ {\lambda} ^ {\alpha (\lambda)}\tag{14}
\]

（乘积取在交换代数 \(\mathbf{S}(\mathbf{M})\) 中）。然后，当 \(\alpha\) 遍历集合 \(\mathbf{N}^{(L)}\) （由从 \(\mathbf{L}\) 到 \(\mathbf{N}\) 的具有有限支撑的映射组成）时，诸 \(e^{\alpha}\) 构成 A-模 \(\mathbf{S}(\mathbf{M})\) 的一组基。

由于 M 是诸 \(M_{\lambda} = Ae_{\lambda}\) 的直和，只需在 L 约化为单个元素时证明该定理，然后应用命题 9 即可。但当 M = Ae（e 为自由元素）时， \(x \otimes y - y \otimes x = 0\) 对所有 x, y 属于 M 成立；理想 \(J'\) 因此为零，从而 \(\mathsf{T}(\mathsf{A}e) = \mathsf{S}(\mathsf{A}e)\) ，于是定理由 § 5，第 5 号，定理 1 得出。

基(14)的乘法表由下式给出

\[
e ^ {\alpha} e ^ {\beta} = e ^ {\alpha + \beta}\tag{15}
\]

（其中 \(\alpha + \beta\) 是从 L 到 N 的映射 \(\lambda \mapsto \alpha(\lambda) + \beta(\lambda)\) 。换言之，S(M) 的基 \((e^{\alpha})\) 连同乘法法则 (15) 典范同构于由 L 导出的自由交换幺半群 \(\mathbf{N}^{(L)}\) ；由此得出（§ 2，第 9 号）：基的指标集为 L 的自由模 M 的对称代数 S(M) 典范同构于多项式代数 A{[}(X \(_{\lambda}\) ) \(_{\lambda \in L}\) {]}，典范同构通过将 \(e_{\lambda}\) 映为 X \(_{\lambda}\) 而得到。特别地（§ 2，第 7 号，命题 7），对于每个从 L 到交换 A-代数 E 中的映射 \(f\colon\) L \(\to\) E ，存在唯一的 A-代数同态 \(\bar{f}\colon\) S(M) \(\to\) E 使得 \(\bar{f}(e_{\lambda}) = f(\lambda)\) 。

注记（2）。上述结果同样可以作为多项式代数和对称代数的泛性质的一个推论而得到，同时考虑到II，§1，第11号，命题17的推论3。

推论。若M是投射A-模，则S(M)是投射A-模。

证明与 § 5，第 5 号，定理 1 的推论的证明相同，只需将 T 换为 S。

命题 10. 设 \(\Delta\) 是一个交换幺半群， \(\mathbf{M}\) 是一个 \(\Delta\) 型分次 A-模， \((\mathbf{M}_{\alpha})_{\alpha \in \Delta}\) 是它的分次。对于每个有序对 \((\alpha, n) \in \Delta \times \mathbf{N}\) ，设 \(\mathbf{S}^{\alpha, n}(\mathbf{M})\) 是 \(\mathbf{S}^n(\mathbf{M})\) 的这样一个子模：诸子模 \(\bigotimes_{\lambda \in J} \mathbf{S}^{n_{\lambda}}(\mathbf{M}_{\alpha_{\lambda}})\) 的直和，其中 \((n_{\lambda})_{\lambda \in L}\) 遍历 \(\geqslant 0\) 的整数族中满足 \(\sum_{\lambda \in L} n_{\lambda} = n, J\) 是其支集的那些族组成的集合，且对每个 \((n_{\lambda})\) , \((\alpha_{\lambda})_{\lambda \in J}\) 遍历 \(\Delta^J\) 的族中满足 \(\sum_{\lambda \in J} \alpha_{\lambda} = \alpha\) 的那些族组成的集合。则 \((\mathbf{S}^{\alpha, n}(\mathbf{M}))_{(\alpha, n) \in \Delta \times \mathbf{N}}\) 是唯一的 \(\Delta \times \mathbf{N}\) 型分次，它与 \(\mathbf{S}(\mathbf{M})\) 上的代数结构相容，并且在 \(\mathbf{M} = \mathbf{S}^1(\mathbf{M})\) 上诱导出给定的分次。

\(\mathbf{S}(\mathbf{M})\) 是诸 \(\mathbf{S}^{\alpha,n}(\mathbf{M})\) 的直和这一事实由命题 9 的推论得出；证明的其余部分与 § 5，第 5 号，命题 7 的证明的结尾完全相同。

更特别地假设 \(\Delta = \mathbf{Z}\) ，并给 \(\mathbb{S}(\mathbf{M})\) 赋予总分次（ \(\mathbf{Z}\) 型）（II，§11，第1号），它对应于上面定义的 \(\mathbf{Z} \times \mathbf{N}\) 型（因而也是 \(\mathbf{Z} \times \mathbf{Z}\) 型）分次；在此分次下次数为 \(n \in \mathbf{Z}\) 的齐次元素因此是属于诸 \(\mathbb{S}^{q, m}(\mathbf{M})\) ——对 \(q + m = n\) ——的直和的那些元素。设 \(f\) 是分次 A-模 \(\mathbf{M}\) 到交换分次 A-代数 \(\mathbf{F}\) （ \(\mathbf{Z}\) 型）的 0 次齐次线性映射；则代数同态 \(g: \mathbb{S}(\mathbf{M}) \to \mathbf{F}\) ——它满足 \(f = g \circ \phi_{\mathbf{M}}'\) ——是一个 \(\mathbf{Z}\) 型分次代数同态，这由公式 \(g(x_1 x_2 \ldots x_n) = f(x_1) f(x_2) \ldots f(x_n)\) （对齐次的 \(x_i\) ，属于 \(\mathbf{M}\) ）、关于 \(f\) 的假设以及 \(\mathbf{Z}\) 型分次在 \(\mathbb{S}(\mathbf{M})\) 上的定义得出。

